\documentclass{article}
\usepackage{amsmath, amssymb, amsthm}
\usepackage{hyperref}
\usepackage{geometry}
\geometry{margin=1in}

\title{Erd\H{o}s Problem \#776}
\date{Last updated: 2025-08-31}
\author{Source: erdosproblems.com}

\begin{document}
\maketitle

\noindent\textbf{Status:} OPEN \quad \textbf{No prize}

\noindent\textbf{Tags:} combinatorics

\medskip
\noindent\textbf{Problem Statement:}

Let $r\geq 2$ and $A_1,\ldots,A_m\subseteq \{1,\ldots,n\}$ be such that $A_i\not\subseteq A_j$ for all $i\neq j$ and for any $t$ if there exists some $i$ with $\lvert A_i\rvert=t$ then there must exist at least $r$ sets of that size.

How large must $n$ be (as a function of $r$) to ensure that there is such a family which achieves $n-3$ distinct sizes of sets?

\bigskip
\noindent\textbf{Remarks:}

A problem of Erd\H{o}s and Trotter. For $r=1$ and $n>3$ the maximum possible is $n-2$. For $r>1$ and $n$ sufficiently large $n-3$ is achievable, but $n-2$ is never achievable.

Let $n_0(r)$ be such that whenever $n>n_0(r)$ there exists such a family with $n-3$ distinct set sizes. He and Tang \cite{HeTa26b} (with ChatGPT) have proved that $n_0(2)=3$, $n_0(3)=8$, and for all $r\geq 4$\[2r+2\leq n_0(r)\leq 2r+2\log_2r+O(\log\log r),\]and hence $n_0(r)\sim 2r$ as $r\to \infty$.

\bigskip
\noindent\textbf{References:}

[HeTa26b] Y. He and Q. Tang, An Erd\H{o}s--Trotter problem on antichains with multiplicity r on each occurring level. arXiv:2602.09803 (2026).

\bigskip
\noindent\small{Source: \url{https://www.erdosproblems.com/776}}
\end{document}
